\section*{Введение}
\addcontentsline{toc}{section}{Введение}
\label{sec:4}
Требуется создать Android/iOS приложение, которым будут пользоваться, как и пациенты, так и врачи.
У первых должна быть возможность проходить опрос с несложными вопросами об их здоровье, у вторых - просмотр результатов и своевременное уведомление врачей о ключевых изменениях.  \par
Каждый пациент должен проходить данный опрос раз в две недели.
Если опрос пройден хорошо, и нет явных признаков ухудшения здоровья, то ему об этом сообщат.
В противном случае, ему должна высветится надпись: \("\)Скоро Вам позвонит врач,\("\)\, а лечащему врачу придет уведомление об ухудшении здоровья его пациента.
Таким образом будет происходить контроль за состоянием здоровья человека. \par
Вредные привычки также играют важную роль в реабилитации.
Поэтому стоит уведомлениями напоминать пациенту о негативных последствиях данных привычек, а также вести контроль за ними.



\subsection*{Общие проблемы нативной разработки iOS-приложений}\label{sec:ios}
Проблемы указаны с учетом специфики проекта, много проблем связанных с финансами, данная платформа не предполагает коммерциализации, поэтому эти проблемы не описаны.

\begin{enumerate}
    \item \textbf{App Store и политика Apple}
    \begin{itemize}
        \item Строгая модерация: долгие проверки, возможные отклонения по формальным причинам.
    \end{itemize}


    \item \textbf{CI/CD и сборка}
    \begin{itemize}
        \item Требуется macOS для сборки (Xcode).
        \item Работа с сертификатами, provisioning profiles -- сложная и неочевидная.
    \end{itemize}

    \item \textbf{Зависимость от Xcode}
    \begin{itemize}
        \item Частые обновления Xcode и macOS.
        \item Нельзя собрать под новую iOS без актуального инструментария.
    \end{itemize}

    \item \textbf{Тестирование}
    \begin{itemize}
        \item Ограничения симулятора: нет камеры, Bluetooth и т.\,д.
        \item TestFlight требует времени и регистрации для внешних тестеров.
    \end{itemize}

\end{enumerate}

\subsection*{Проблемы разработки iOS-приложений в России}

\begin{enumerate}
    \item \textbf{Невозможность прямой публикации из РФ}
    \begin{itemize}
        \item Apple не принимает карты российских банков.
        \item Регистрация разработчика из РФ невозможна.
        \item Обход: регистрация через третьи страны (Казахстан, Грузия, Армения и т.\,д.).
    \end{itemize}

    \item \textbf{Риски блокировок}
    \begin{itemize}
        \item Приложения российских банков, медиа и телекомов уже удалялись.
        \item Возможна блокировка аккаунта и удаление приложения без предупреждения.
    \end{itemize}

    \item \textbf{Отсутствие монетизации в РФ}
    \begin{itemize}
        \item App Store не работает полноценно в России.
        \item Большая часть аудитории не может платить через стандартные механизмы.
    \end{itemize}

\end{enumerate}